% ===========================================================================
%  E14 — Evidência pedagógica
% ===========================================================================
\chapter{E14 \textperiodcentered{} Evidência pedagógica}
\label{ep:e14}

\begin{fichaepico}
\begin{campos}
  \item[Objetivo] A plataforma registra, para cada estudante, quais habilidades
    da BNCC ele demonstrou numa atividade de uma missão, rodada a rodada, e o
    professor lê, corrige e leva esse retrato para a avaliação dele.
  \item[Para quem] Professor (lê, confirma e ajusta) e Estudante (responde às
    perguntas da missão, sem ver o próprio resultado). A Coordenação e a
    Administração consultam, sem mexer. A equipe pedagógica da Unooria valida a
    régua de cada missão.
  \item[Valor] A missão já via o que a criança fazia --- tirou o post do ar,
    respondeu a pergunta ---, e isso ficava dentro dela. No fim do bimestre, o
    professor não tinha como dizer \enquote{a Ana atingiu a EF08CO05 e o Caio
    ainda está em processo}, que é a frase que a rede pede e que a BNCC
    organiza.
  \item[Estado] \parcial{} --- a base está pronta na plataforma; aparece para o
    professor quando a primeira missão tiver a régua pedagógica, e nenhuma tem
    ainda.
\end{campos}
\end{fichaepico}

\begin{tabularx}{\linewidth}{@{}L{16mm}X L{20mm}@{}}
\toprule
\headrow \thd{ID} & \thd{Feature} & \thd{Estado}\\
\midrule
\id{F14.1} & A régua pedagógica de cada missão & \parcial\\
\zebra \id{F14.2} & O registro por atividade e rodada & \parcial\\
\id{F14.3} & Equipes e aula projetada & \parcial\\
\zebra \id{F14.4} & O painel e o ajuste do professor & \parcial\\
\id{F14.5} & A pergunta de múltipla escolha & \parcial\\
\zebra \id{F14.6} & O perfil do aluno & \parcial\\
\bottomrule
\end{tabularx}

\section{O que é uma evidência aqui}

\textbf{Evidência é o nível que um estudante demonstrou numa habilidade da
BNCC}, dentro de uma atividade de uma missão, numa rodada. Três ideias a
sustentam:

\begin{itemize}
  \item \textbf{A atividade é da missão.} Cada missão diz quais são as suas
    atividades; a plataforma não tem, por ora, uma \enquote{atividade} dela
    mesma. Hoje só o Instanooria tem atividades --- a que o professor cria a
    partir da aula (\id{H06.6.1}) ---, e o unooGPT e o UnooTube serão refeitos
    com atividades numa próxima onda.
  \item \textbf{A rodada é cada vez que a atividade é feita.} Reiniciar a
    atividade abre uma rodada nova, e a anterior fica guardada.
  \item \textbf{A régua pedagógica diz o que a missão mede}: quais habilidades,
    em quais momentos, qual pergunta e quanto vale cada alternativa. Ela é
    escrita com a equipe pedagógica e validada por ela antes de entrar.
\end{itemize}

A escala tem três níveis --- \textbf{atingiu}, \textbf{em processo} e
\textbf{não atingiu} ---, e a falta de registro é \enquote{Não avaliado}, que
nunca conta como não atingiu.

\begin{nota}
\textbf{Pronta na plataforma, invisível para o professor.} As histórias
marcadas \parcial{} neste épico têm a regra construída e testada, mas nenhuma
missão foi adaptada: sem régua, a página da atividade não mostra o bloco
\enquote{Habilidades da BNCC}, e o perfil do estudante fica vazio. Elas passam a
\pronto{} quando a primeira missão entrar com a régua (\id{H14.1.2}) e houver o
teste que percorre a tela.
\end{nota}

\section{O que foi decidido}

\begin{decisao}
\textbf{Decidido em 25 de setembro de 2026.}
\begin{itemize}
  \item \textbf{Esta frente é só a plataforma.} Nenhuma missão foi adaptada; cada
    uma entra depois, com a régua validada pela equipe pedagógica.
  \item \textbf{Três níveis e um \enquote{Não avaliado}.} Sem registro, a
    habilidade é \enquote{Não avaliado}, e nunca conta como não atingiu.
    \enquote{Não participou} é o \enquote{Não avaliado} que o professor põe à
    mão.
  \item \textbf{Vale o melhor.} Dentro de uma rodada, dois momentos que medem a
    mesma habilidade ficam com o melhor; o resultado da atividade é a melhor
    rodada de cada habilidade.
  \item \textbf{Quem vale mais}: o professor, depois o próprio estudante, depois
    a equipe. A equipe recebe só o nível, nunca texto nenhum.
  \item \textbf{Trocar a equipe muda só a rodada em andamento.} As rodadas
    passadas ficam com a equipe que tinham, e o que estiver errado nelas o
    professor corrige à mão.
  \item \textbf{Diferente do documento do PO: a aula projetada não vale
    sozinha.} O pedido era replicar automaticamente para todos, também quando a
    turma inteira usa um computador só. Na projeção, o resultado chega \enquote{a
    revisar} e só conta depois de \enquote{Confirmar para a turma}: a criança no
    fundo da sala não demonstrou o que a turma respondeu no telão.
  \item \textbf{Cada missão tem as suas atividades.} Uma atividade comum a todas
    as missões fica para quando houver mais de uma missão com atividades.
  \item \textbf{A resposta certa fica na plataforma.} Na pergunta de múltipla
    escolha vale a primeira resposta, e o estudante só vê \enquote{resposta
    enviada}.
  \item \textbf{O estudante não vê o próprio resultado, e não há ranking.}
  \item \textbf{O perfil é por missão e por atividade.} A mesma habilidade
    medida em duas missões aparece nas duas; juntar as missões é de um painel
    futuro, com regra pedagógica própria (\id{H14.6.2}).
\end{itemize}
\end{decisao}

% ---------------------------------------------------------------------------
\feature{F14.1}{A régua pedagógica de cada missão}{\parcial}

\begin{historia}{H14.1.1}{Cada missão diz o que mede}{\parcial}
\hu{Equipe pedagógica da Unooria}{definir, para cada missão, quais habilidades
da BNCC ela mede e em quais momentos}{que o painel do professor afirme só o que
a missão de fato observa}
\aceite
\begin{itemize}
  \item A régua lista as habilidades pelo código exato da BNCC (ex.:
    EF08CO05), os momentos que medem cada uma --- uma ação dentro da missão ou
    uma pergunta de múltipla escolha --- e, em cada pergunta, quanto vale cada
    alternativa: atingiu, em processo ou não atingiu.
  \item A régua diz quem a validou na equipe pedagógica, e quando.
  \item A plataforma não aceita régua com habilidade que não existe na BNCC,
    com alternativa fora da escala ou com pergunta que nenhum momento corrige.
  \item Mudou a régua, muda a versão; cada registro guarda a versão que valia
    quando foi feito.
  \item Missão sem régua não mostra evidência nenhuma.
\end{itemize}
\pendencia{a base está pronta; nenhuma missão tem régua ainda (\id{H14.1.2}).}
\end{historia}

\begin{historia}{H14.1.2}{A primeira missão entra com a régua}{\afazer}
\hu{Professor}{que a primeira missão com atividades venha com a régua
pedagógica}{ver, na página da atividade, as habilidades que cada estudante
demonstrou}
\aceiteproposta
\begin{itemize}
  \item A régua é escrita com a equipe pedagógica e validada antes de entrar.
  \item A missão registra cada momento da régua enquanto a turma joga.
  \item O bloco \enquote{Habilidades da BNCC} aparece na página da atividade.
  \item A primeira candidata é o Instanooria, a única missão que já tem
    atividades.
\end{itemize}
\end{historia}

% ---------------------------------------------------------------------------
\feature{F14.2}{O registro por atividade e rodada}{\parcial}

\begin{historia}{H14.2.1}{Um nível por habilidade, e \enquote{Não avaliado} quando não há}{\parcial}
\hu{Professor}{ver, para cada estudante e cada habilidade, o nível que ele
demonstrou}{saber quem atingiu e quem ainda está em processo}
\aceite
\begin{itemize}
  \item A escala tem três níveis: atingiu, em processo e não atingiu.
  \item Sem registro, a habilidade aparece como \enquote{Não avaliado} --- e
    nunca conta como não atingiu.
  \item Quando dois momentos da mesma rodada medem a mesma habilidade, vale o
    melhor: demonstrar uma vez é demonstrar.
\end{itemize}
\pendencia{a base está pronta; aparece para o professor quando a primeira
missão tiver a régua pedagógica.}
\end{historia}

\begin{historia}{H14.2.2}{Refazer a atividade sem perder o que já foi demonstrado}{\parcial}
\hu{Professor}{reiniciar uma atividade para a turma tentar de novo}{dar outra
chance sem apagar o que cada estudante já tinha mostrado}
\aceite
\begin{itemize}
  \item Cada vez que a atividade é feita é uma rodada: reiniciar abre a rodada
    seguinte, e a anterior fica guardada.
  \item O resultado da atividade é a melhor rodada de cada habilidade: quem
    atingiu uma habilidade na primeira rodada e outra na segunda atingiu as
    duas.
  \item Uma rodada pior depois não apaga o que já tinha sido demonstrado.
\end{itemize}
\pendencia{a base está pronta; aparece para o professor quando a primeira
missão tiver a régua pedagógica.}
\end{historia}

\begin{historia}{H14.2.3}{Sem ranking, e o nível é da habilidade}{\parcial}
\hu{Professor}{ler a evidência da turma sem ranking nem nota somada}{avaliar
cada habilidade sem transformar a criança num número}
\aceite
\begin{itemize}
  \item Os estudantes aparecem em ordem alfabética, nunca por desempenho.
  \item Nenhuma tela soma os níveis numa nota: o nível é de uma habilidade, não
    da criança.
\end{itemize}
\pendencia{a base está pronta; aparece para o professor quando a primeira
missão tiver a régua pedagógica.}
\end{historia}

% ---------------------------------------------------------------------------
\feature{F14.3}{Equipes e aula projetada}{\parcial}

Quando os estudantes dividem um computador (\id{E13}), só um deles entra na
missão. O que a missão vê naquela tela foi feito pela equipe.

\begin{historia}{H14.3.1}{Quem divide o computador recebe o mesmo nível}{\parcial}
\hu{Professor}{que o que a equipe fez junto valha para todos da equipe}{não ver
como \enquote{Não avaliado} quem estava sentado ao lado de quem entrou}
\aceite
\begin{itemize}
  \item Quando o estudante que entrou demonstra uma habilidade, os colegas do
    mesmo computador recebem o mesmo nível.
  \item Vai só o nível: nada do que foi escrito passa de um para o outro.
  \item Quem entrou com o próprio login e passou por aquele momento fica com o
    que ele mesmo fez, ainda que o colega tenha ido melhor.
  \item O que o professor ajustou nunca é trocado pelo que veio da equipe.
  \item Se dois da mesma equipe demonstraram a habilidade, fica o melhor.
\end{itemize}
\pendencia{a base está pronta; aparece para o professor quando a primeira
missão tiver a régua pedagógica.}
\end{historia}

\begin{historia}{H14.3.2}{Trocar as equipes no meio da atividade}{\parcial}
\hu{Professor}{mudar as equipes depois que a atividade começou}{corrigir quem
sentou onde sem mexer nas aulas anteriores}
\aceite
\begin{itemize}
  \item Ao salvar ou desfazer as equipes, a plataforma refaz o que veio da
    equipe só na rodada em andamento.
  \item As rodadas passadas ficam com a equipe que tinham naquela aula; o que
    estiver errado nelas eu corrijo à mão (\id{H14.4.2}).
  \item Quem saiu da equipe perde o nível que tinha recebido dela; quem entrou
    recebe.
  \item Salvar as equipes de novo sem mudar ninguém de lugar não desfaz a
    confirmação que eu já tinha dado (\id{H14.3.3}).
\end{itemize}
\pendencia{a base está pronta; aparece para o professor quando a primeira
missão tiver a régua pedagógica.}
\end{historia}

\begin{historia}{H14.3.3}{A aula projetada chega \enquote{a revisar}}{\parcial}
\hu{Professor}{que o resultado da turma inteira diante de um computador só não
valha para todos sem eu conferir}{não afirmar que cada criança demonstrou o que
a turma respondeu junto}
\aceite
\begin{itemize}
  \item Quando a turma inteira usa um computador --- as equipes dizem que há um
    computador, ou uma equipe tem todos os estudantes ---, o nível chega a cada
    um como \enquote{a revisar}.
  \item O que está \enquote{a revisar} não conta no resultado da atividade nem
    no perfil do estudante.
  \item \enquote{Confirmar para a turma} faz valer, de uma vez, tudo o que está
    a revisar naquela rodada; confirmar de novo não muda nada.
  \item Depois de confirmar, ainda posso ajustar cada estudante (\id{H14.4.2}).
  \item É diferente do documento do PO, que pedia a replicação automática
    também aqui.
\end{itemize}
\pendencia{a base está pronta; aparece para o professor quando a primeira
missão tiver a régua pedagógica.}
\end{historia}

% ---------------------------------------------------------------------------
\feature{F14.4}{O painel e o ajuste do professor}{\parcial}

\begin{historia}{H14.4.1}{Ver as habilidades da turma na página da atividade}{\parcial}
\hu{Professor}{ver, na página da atividade, as habilidades da BNCC de cada
estudante}{saber o que a turma demonstrou e planejar a próxima aula}
\aceite
\begin{itemize}
  \item A página da atividade ganha o bloco \enquote{Habilidades da BNCC} só
    quando a missão tem régua.
  \item Cada estudante da turma é uma linha, com a equipe dele na missão, o
    nível de cada habilidade em cada rodada e o resultado da atividade.
  \item Cada habilidade vem com o texto da BNCC, e cada nível diz de onde veio:
    do próprio estudante, da equipe, da aula projetada ou do ajuste do
    professor.
  \item O bloco diz quantos níveis de cada rodada estão \enquote{a revisar}.
  \item Hoje só o Instanooria tem atividades, então é só na página da
    atividade dele que o bloco pode aparecer.
\end{itemize}
\pendencia{a base está pronta; aparece para o professor quando a primeira
missão tiver a régua pedagógica.}
\end{historia}

\begin{historia}{H14.4.2}{Ajustar um estudante ou vários}{\parcial}
\hu{Professor}{corrigir o nível de um estudante, ou de vários de uma vez}{fazer
valer o que eu vi na sala, mesmo quando a missão não viu}
\aceite
\begin{itemize}
  \item Posso ajustar a qualquer momento, mesmo com a missão encerrada.
  \item Posso ajustar um estudante ou vários de uma vez, numa habilidade e numa
    rodada.
  \item O ajuste vale sobre tudo: nem o que o estudante fizer depois, nem a
    equipe, o trocam.
  \item Cada mudança fica no histórico --- de quê para quê, quem mudou e quando
    ---, e o histórico não pode ser reescrito.
  \item Só quem dá aula na turma ajusta.
\end{itemize}
\pendencia{a base está pronta; aparece para o professor quando a primeira
missão tiver a régua pedagógica.}
\end{historia}

\begin{historia}{H14.4.3}{Marcar que o estudante não participou}{\parcial}
\hu{Professor}{marcar que um estudante não participou de uma atividade}{que
quem faltou não herde o que a equipe fez}
\aceite
\begin{itemize}
  \item \enquote{Não participou} é um \enquote{Não avaliado} que eu mesmo
    ponho.
  \item Ele impede que o nível da equipe chegue a esse estudante, também quando
    as equipes mudam.
  \item No perfil, conta como não avaliado --- nunca como não atingiu.
  \item Só o professor marca; a missão nunca marca \enquote{não participou}.
\end{itemize}
\pendencia{a base está pronta; aparece para o professor quando a primeira
missão tiver a régua pedagógica.}
\end{historia}

\begin{historia}{H14.4.4}{A coordenação consulta, sem ajustar}{\parcial}
\hu{Coordenação da escola}{consultar a evidência de uma turma e o perfil de um
estudante}{acompanhar a aprendizagem sem interferir na avaliação do professor}
\aceite
\begin{itemize}
  \item A Coordenação da escola, a Coordenação regional, o Administrador do
    contrato e o Administrador Unooria podem ler a evidência de uma turma e o
    perfil de um estudante.
  \item Não há confirmar nem ajustar para eles: quem decide é quem dá aula na
    turma.
  \item O que se lê é o nível de cada habilidade, nunca o que a criança
    escreveu.
  \item É a exceção declarada à regra de que a gestão não lê a linha de uma
    criança (\id{H07.1.3}): aqui ela lê o nível por habilidade, e nunca o
    relatório da missão.
\end{itemize}
\pendencia{a plataforma já permite a leitura e já recusa o ajuste da
coordenação, mas as telas desta frente são as do professor: a tela da
coordenação ainda não existe.}
\end{historia}

% ---------------------------------------------------------------------------
\feature{F14.5}{A pergunta de múltipla escolha}{\parcial}

\begin{historia}{H14.5.1}{Responder sem saber se acertei}{\parcial}
\hu{Estudante}{responder à pergunta da missão}{mostrar o que eu sei sem que a
missão vire uma prova com nota na hora}
\aceite
\begin{itemize}
  \item A pergunta mostra o enunciado e as alternativas; quanto vale cada
    alternativa fica na plataforma, e nunca chega ao meu aparelho.
  \item Depois de responder, vejo só \enquote{resposta enviada} --- e vejo o
    mesmo se responder de novo.
  \item Vale a primeira resposta; a segunda não muda nada.
  \item A resposta conta na atividade e na rodada em que eu estou na missão:
    não sou eu quem escolhe.
  \item Se a missão não está aberta na minha turma, a pergunta não abre.
  \item Não vejo o meu resultado em lugar nenhum da plataforma.
\end{itemize}
\pendencia{a base está pronta; aparece quando a primeira missão tiver a régua
pedagógica, com as perguntas dela.}
\end{historia}

% ---------------------------------------------------------------------------
\feature{F14.6}{O perfil do aluno}{\parcial}

\begin{historia}{H14.6.1}{Abrir o perfil de um estudante}{\parcial}
\hu{Professor}{abrir o perfil de um estudante da minha turma}{ver, missão a
missão, as habilidades que ele já demonstrou}
\aceite
\begin{itemize}
  \item O perfil abre pela aba Estudantes da turma.
  \item Em cada missão e em cada atividade, cada habilidade aparece com o nível
    da melhor rodada.
  \item Entra só o que vale: nada \enquote{a revisar}, nada de \enquote{não
    participou}, nada de \enquote{Não avaliado}.
  \item O perfil diz o que o estudante demonstrou, e não de onde veio (dele, da
    equipe ou do professor).
  \item A mesma habilidade medida em duas missões aparece nas duas, separada.
  \item Só abro o perfil de quem eu ensino; outro estudante aparece como não
    encontrado.
\end{itemize}
\pendencia{a base está pronta; o perfil fica vazio até a primeira missão ter a
régua pedagógica.}
\end{historia}

\begin{historia}{H14.6.2}{Juntar as habilidades de todas as missões}{\afazer}
\hu{Professor}{ver num painel só cada habilidade da BNCC de um estudante,
somando todas as missões}{acompanhar a Base ao longo do ano, e não missão por
missão}
\aceiteproposta
\begin{itemize}
  \item A regra de junção --- vale a melhor? a mais recente? pesa o tamanho da
    missão? --- é definida pela equipe pedagógica antes de qualquer tela.
  \item O painel mostra de quais missões veio cada nível.
  \item Continua sem ranking e sem nota somada.
\end{itemize}
\end{historia}
